Beaglehole et al. \cite{beaglehole2026toward} claim that the concepts ``known'' by Large Language Models are in fact linear directions in the model's activation space, that these directions can be extracted with feature-learning algorithms such as the Recursive Feature Machines (hereafter RFM), and that, once a concept has been extracted, one may \textbf{steer} the model, that is, push its outputs toward or away from that concept by purely linear computations, or \textbf{monitor} the presence of the concept by reading the model's activations. These claims are established over Llama-class models with 8 to 90 billion parameters, and it is therefore reasonable to ask whether the same results survive at much smaller scales. Do such findings endure on a model as small as $\approx 9$ million parameters? Note that a single checkpoint cannot say how the machinery scales, since that would require a ladder of model sizes and is a separate experiment altogether, but it can say whether the mechanism survives at the bottom of the range on its own merits, and a negative answer there would bound the whole claim.

This document reports a controlled replication of the aforementioned mechanisms, namely RFM extraction, additive steering, AGOP-eigenvector probing and the comparison against cheaper baselines, carried out on GuppyLM, a language model of $\approx 8.7$M parameters. We deliberately dropped the themed results of Beaglehole et al., such as steering along an anti-refusal direction to expose a model's jailbreak vulnerability, or subtracting a deception direction so that the model answers more honestly, since GuppyLM does not carry representations of such matters. Instead we ask whether linear concept extraction, additive steering and activation probing hold on language models as small as 10M parameters, restricted to the concepts the model is able to represent.

In short, the reading of concepts transfers intact, the writing of them, that is, the steering, transfers as a mechanism and yet not as a capability, and the baseline comparison is only partly testable at this scale.
